\part{Rahmen- bedingungen für KI-Sicherheit in der Luftfahrt}
\chapterimage{Bilder_und_Co/safety_sign_with_computer_code_and_a_airplane_in_the_background_1958542794.png} % Chapter heading image
%----------------------------------------------------------------------------------------
\chapter{Luftfahrt}

% ==== Luftfahrt ====
\subsection{Ziel, Motivation und Scope (Luftfahrt)}
\paragraph{Ziel.}
Einordnung, wie „sicher genug“ für KI-Funktionen in der Luftfahrt regulatorisch begründet und nachgewiesen wird: Rolle von EASA (AI Roadmap, Concept Paper Issue~2), MLEAP-Ergebnissen und etablierten Standards (z.B. DO-178C). Klarstellung: KI-Guidance ergänzt, ersetzt aber nicht die klassische Assurance; keine Erleichterung der Nachweisstrenge.

\paragraph{Motivation.}
KI hält in sicherheitsrelevanten Avionikfunktionen Einzug (z.\,B.\ Radar-Objektklassifizierung). Für den Einsatz im Flugbetrieb braucht es eine belastbare Begründung „sicher genug“ – auf Basis der etablierten Luftfahrt-Standards, ergänzt um KI-spezifische Evidenzen. Gleichzeitig bleiben zentrale Fragen offen (Generalisierung und Drift, Daten-/Modell-Governance, Updates im Betrieb, OOD/Abstain). Ein kurzer Seitenblick auf militärische Anwendungen ist hilfreich, weil dort Degradation, Resilienz gegen Störungen (Jamming/Täuschung) und konsequente Human-in-the-Loop-Konzepte geschärft sind; diese Aspekte liefern Anregungen, stehen hier aber nicht im Mittelpunkt der zivilen Zulassung.

\paragraph{Scope.}
Fokus auf Safety-Aspekte von KI in der Avionik, insbesondere \emph{Radar-Objektklassifizierung} als Beispielanwendung. Abdeckung der EASA-Guidance zu Level~1/2 (Assistenz, Human-AI-Teaming), \emph{offline} trainierten, zur Freigabe eingefrorenen Modellen.





\section{EASA AI-Roadmap – aktueller Stand}
Die EASA positioniert KI als „human-centric“ und staffelt Anwendungen nach Autorität: \textbf{Level 1} (Assistenz), 
\textbf{Level 2} (Human-AI-Teaming) und später \textbf{Level 3} (fortgeschrittene Automatisierung). 
Roadmap 2.0 definiert Vertrauenswürdigkeits-Bausteine (u.\,a. Learning Assurance, Erklärbarkeit, Human Oversight) 
und skizziert einen gestuften Zulassungspfad. \emph{Kernbotschaft:} KI ergänzt, aber ersetzt das bestehende Safety-Regelwerk nicht.
\paragraph{Was das praktisch bedeutet.}
\begin{itemize}
  \item \textbf{Stufenweise Zulassung:} Erst Level-1, später Level-2 – jeweils mit klar definiertem Einsatzbereich und menschlicher Aufsicht. \emph{Kein „Assurance-Rabatt“:} Die Strenge der Nachweise richtet sich weiterhin nach Hazard/DAL.
  \item \textbf{Zusatz statt Ersatz:} KI-spezifische Evidenzen (z.\,B. Daten-/Modell-Governance, Kalibrierung, OOD/Abstain, Monitoring) kommen \emph{on top} zur etablierten Software-/System-Assurance.
  \item \textbf{Lebenszyklus im Blick:} Betriebliches Monitoring, kontrollierte Updates und nachvollziehbare Änderungen sind Teil des Sicherheitsarguments („living“ Safety Case).
\end{itemize}


\section{EASA AI-Guidance 2024 (Concept Paper Issue 2) – Level 1 \& 2}
Das \emph{Concept Paper Issue 2} (03/2024) beschreibt die schrittweise Einführung von KI in sicherheitsrelevanten Anwendungen: zuerst \textbf{Level 1} (Assistenzfunktionen), anschließend ausgewählte \textbf{Level 2}-Anwendungen (\emph{Human–AI Teaming} unter klarer menschlicher Aufsicht). Die Guidance ergänzt das bestehende Safety-Regelwerk (z.\,B.\ DO-178C) um KI-spezifische Bausteine wie \emph{Learning Assurance}, Erklärbarkeit, präzise Einsatzgrenzen (ODD) und betriebliches Monitoring – ohne die klassischen Nachweise zu ersetzen.

\paragraph{Schwerpunkte.}
\begin{itemize}
  \item \textbf{Klare Rollenverteilung:} L1 = Unterstützung; L2 = Kooperation mit dem Menschen. Autorität, Eingriffsrechte und Übersteuerbarkeit müssen beschrieben und gegeben sein.
  \item \textbf{Einsatzbereich festlegen (ODD):} Scope, Use Cases, Datenbasis und Grenzen definieren – inkl.\ Regeln, wann „Abstain/Unknown“ zu ziehen ist.
  \item \textbf{Learning Assurance:} Herkunft/Qualität der Trainingsdaten, kontrollierte Modelländerungen und Freigabeprozesse transparent darlegen.
  \item \textbf{Erklärbarkeit \& HMI:} Begründbare Ausgaben, Konfidenz-/Qualitätsanzeigen und eindeutige Degradationsmodi.
  \item \textbf{Betrieb \& Monitoring:} Drift-Überwachung, Logging, Eskalationswege; \emph{„living“ Safety Case} statt Einmal-Nachweis.
\end{itemize}

\paragraph{Frühe Einsatzgrenzen.}
\begin{itemize}
  \item \textbf{Begrenzte Autorität:} L1 = Mensch entscheidet und handelt; L2 = System darf handeln, bleibt aber jederzeit übersteuerbar.
  \item \textbf{Kein Online-Lernen:} Zulässig sind \emph{offline} gelernte, zur Freigabe \emph{eingefrorene} Modelle; Online-/Continual Learning ist (noch) ausgeschlossen.
  \item \textbf{ODD \& OOD:} Der Einsatzbereich ist eng zu beschreiben; für \emph{außerhalb der ODD}/\emph{Out-of-Distribution} sind Monitoring, Warnung und sicheres Verhalten gefordert.
  \item \textbf{Nachweisstrenge bleibt:} KI-spezifische Evidenzen kommen \emph{on top}; die Software-/Hardware-Assurance (z.\,B.\ DO-178C/DO-254) bleibt vollumfänglich anwendbar – es gibt keine Erleichterung der Nachweisstrenge.
  \item \textbf{Zielniveau:} Mindestens das Sicherheitsniveau konventioneller Lösungen.
\end{itemize}




\section{MLEAP Final Report (2024) – Kernaussagen und Relevanz}
Der \emph{Final Report} des EASA-Forschungsprojekts \textbf{MLEAP} (Machine Learning Application Approval) bündelt Methoden- und Werkzeugempfehlungen zur \emph{Learning Assurance} und zeigt, wie der EASA-\emph{W-Shape} (lernprozessspezifische Ergänzung zum V-Modell) praktisch umgesetzt werden kann. Schwerpunkte sind \textbf{Datenrepräsentativität}, \textbf{Generalisierung} und \textbf{Robustheit/Stabilität} von ML-Modellen; ergänzt wird dies durch eine generische Entwicklungs-Pipeline im Sinne des W-Shape. Die Ergebnisse fließen in das \emph{EASA AI Concept Paper Issue 2} ein.

\paragraph{Kerninhalte.}
\begin{itemize}
  \item \textbf{W-Shape als Rahmen:} Abgrenzung System-Ebene vs.\ KI-Ebene; spezifische Ziele und Verifikationen für Datenmanagement, Lernprozess, Modellverifikation und Integration werden entlang des Lernzyklus verankert. 
  \item \textbf{Drei technische Achsen:} (1) \emph{Datenvollständigkeit \& Repräsentativität} inkl.\ Corner-/Edge-Cases und Auswahlraster für Datenqualifikationsmethoden; (2) \emph{Generalisierung} (Vermeidung von Unter/Überanpassung, Abschätzung/Bounds, prozessnahe Maßnahmen); (3) \emph{Robustheit \& Stabilität} gegenüber Störungen/Rauschen inkl.\ Ansätzen von empirischen bis formalen Verfahren. 
  \item \textbf{Generische Pipeline:} Ein praktikabler Ablauf zur Umsetzung des W-Shape mit klaren Prüfbausteinen (Daten-, Lern-, Performance- und Integrationsverifikationen) wird vorgeschlagen. 
  \item \textbf{Use-Cases (Evidenz):} \emph{ATC-Speech-to-Text}, \emph{Automated Visual Inspection} (In-Service-Schadenerkennung) und \emph{ACAS Xu} (UAS-Air-to-Air-Avoidance) wurden untersucht, um Methoden und Empfehlungen an realen Luftfahrtaufgaben zu validieren. 
\end{itemize}





\section{DO-178C – Rolle im KI-Kontext}
\textbf{DO-178C} bleibt die \emph{Baseline} für Software-Assurance in der Avionik: Ziele und Evidenzen skalieren mit 
\textbf{DAL A–E}. EASA verweist explizit auf die DO-178C-Daten/Liefergegenstände (inkl. Tool-Qualifikation) als Nachweisgrundlage; 
KI-spezifische Guidance kommt \emph{zusätzlich} obenauf (z.\,B. Learning Assurance), ändert aber die Software-Pflichten nicht. 
\emph{Praxis:} Daten-/Modell-Pipelines berühren Tool-Qualifikation (DO-330).

\paragraph{Bezug zu KI Anwendungen und EASA}
\begin{itemize}
  \item \textbf{Airborne Software bleibt DO-178C-pflichtig:} Inferenz-Software, Laufzeit/Framework, Treiber und Integrationscode (inkl.\ Vor-/Nachverarbeitung) werden wie gewohnt gemäß DO-178C nachgewiesen; die KI-Guidance kommt \emph{on top} (kein „DAL-Rabatt“).
  \item \textbf{Modelle sind konfigurationsgeführte Artefakte:} Gewichte/Graph/Quantisierung gelten als freigegebener Software-Stand („gefroren“ zur Zulassung) und werden wie Software-Daten behandelt; Updates sind kontrollierte Änderungen, kein Online-Lernen.
  \item \textbf{Tool-Qualifikation denken:} Trainings-/Konvertierungs-/Code-Gen-Tools können qualifikationspflichtig sein (DO-330), wenn ihre Fehler unentdeckt ins Produkt wirken könnten.
  \item \textbf{Traceability bleibt zentral:} Anforderungen \(\rightarrow\) Architektur/Code \(\rightarrow\) Tests/Analysen müssen lückenlos verknüpft sein; KI-spezifische Evidenz (z.\,B. Daten-/Modell-Governance, Kalibrierung, OOD/Abstain) wird \emph{zusätzlich} dokumentiert.
\end{itemize}
\paragraph{Wichtig:}
Die DO 178C und 330 nehmen kein Bezug zu KI und ist für Software und Tools im Allgemeinen. Anzumerken ist, dass deterministischer code 
oder MC/DC Coverage stark empfohlen wird und als best practise gilt. Diese Punkte sind bei KI nicht erfüllbar und werden nicht in dieser 
Norm ausdrücklich gefordert. Das Fazit dazu ist, dass KI nach diesen Standards möglich ist, aber die Anforderungen an Nachweisen steigen.






\section{Lückenanalyse: Zertifizierungsfähigkeit von KI in der Luftfahrt}
\textit{``ML meets aerospace: challenges of certifying airborne AI''} hat in den Überlegungen und Limitierungen einige Punkte heraus gedeutet:
\subsection{Überlegungen:}
\paragraph{(1) Geltung heutiger Prozesse für KI (Anwendbarkeit)}
\textbf{Befund:} KI ist formal „nur“ eine Implementierungswahl, verhält sich aber probabilistisch und datenverteiltungsabhängig.  
\textbf{Gap:} Klassische Nachweise decken „eingefrorene“ Leistungswerte ab, nicht jedoch das Verhalten auf \emph{unbekannten} Daten und unter \emph{Drift}. Der ODD-Begriff allein reicht nicht, wenn die reale Datenverteilung wandert.

\paragraph{(2) Safety-Aspekte}
\textbf{Befund:} Anforderungen können Zielwerte (z.\,B. Treffer-/Fehlerraten) vorgeben; Verifikation gelingt auf Validationsdaten.  
\textbf{Gaps:} 
\begin{itemize}
  \item Übertragbarkeit auf \emph{unbekannte} Daten bleibt unsicher (Generalisierung/Drift).
  \item Zwei Wege werden diskutiert: (1) wie bisher über Requirements \& Design Assurance \emph{ohne} probabilistische Systemintegration; (2) \emph{mit} expliziter Integration von Zuverlässigkeitswerten in Safety-Analysen – dafür fehlen heute belastbare Verfahren.
  \item „DALgebra“ (Redundanz/Monitoring \(\rightarrow\) höherer DAL) ist für ML nur eingeschränkt übertragbar: Unabhängigkeit und verlässliche Erkennung „nicht-leistender“ Kanäle sind ungeklärt.
\end{itemize}

\paragraph{(3) System-Aspekte}
\textbf{Befund:} Luftfahrtsysteme nutzen Redundanz und Dissimilarität zur Verringerung gemeinsamer Fehlerursachen.  
\textbf{Gaps:}
\begin{itemize}
  \item \textit{Dissimilarität bei ML}: Was bedeutet „hinreichend verschieden“ bei Daten, Architektur, Hyperparametern, Initialisierung?
  \item \textit{Runtime Assurance}: Umschalten auf einfachere Pfade ist richtig, aber \emph{Erkennung} von Fehlfunktionen, verlässliche Konfidenzen und OOD-Detektion sind weiterhin Forschungsfelder.
\end{itemize}

\paragraph{(4) Software-Aspekte (DO-178C-Logik)}
\textbf{Befund:} Klassische Ebenen (System-, High-, Low-Level-Requirements), Traceability, Coverage und DAL-skalierte Evidenz.  
\textbf{Gaps:}
\begin{itemize}
  \item \textit{Spezifikation}: ML-Verhalten wird \emph{indirekt} über Daten/Metriken beschrieben – wie wird diese „Zwischen-Spezifikation“ standardkonform geführt?
  \item \textit{Traceability}: Vom High-Level-Requirement bis hinunter zu \emph{Datenanforderungen}, Datensätzen und Metriken – Lücke zwischen Spezifikation und tatsächlichem Modellverhalten.
  \item \textit{Verifikation}: Nominal/off-nominal Testdateneigenschaften, Robustheit/Generalisierung – \emph{ausreichende Testtiefe} und \emph{Coverage-Maße} für ML sind offen.
\end{itemize}

\paragraph{(5) Probabilistische Systeme (Analogie)}
\textbf{Befund:} GPS/ILS zeigen, wie probabilistische Leistungsanforderungen regulativ handhabbar sind (Schwellen, Tests, Konfidenzen).  
\textbf{Gap:} Vergleichbare, anerkannte Verfahren fehlen für ML: Wann ist Testen „genug“, und wie wird probabilistische Leistung sicher argumentiert (inkl. OoD/Drift)?

\paragraph{(6) Zertifizierungsprozess}
\textbf{Befund:} Behörden geben Regeln vor; akzeptierte \emph{Means of Compliance} kommen aus Konsensstandards.  
\textbf{Gaps:}
\begin{itemize}
  \item „Henne-Ei“: Reife Praxis entsteht erst im Betrieb, Flugbetrieb setzt aber Zulassung voraus.
  \item Standardisierungszyklen vs.\ ML-Innovationsgeschwindigkeit: Risiko der Obsoleszenz.
  \item Kontinuierliche Updates sind für ML essenziell, triggern aber (heute) Änderung/Zertifizierung – Guidance für \emph{kontinuierliche Re-Zertifizierung} fehlt.
\end{itemize}

\paragraph{(7) Tool-Qualifikation (DO-330)}
\textbf{Befund:} TQL 1–5 je nach Toolkriterium und Ziel-DAL; Qualifikation nötig, wenn Prozesse ersetzt/reduziert werden.  
\textbf{Gaps:}
\begin{itemize}
  \item ML-Toolchains (Datenmanagement, Training, Konvertierung, Deployment, OoD-Checks) sprengen die klassische Tool-Taxonomie.
  \item „Output verifizieren statt Tool qualifizieren“ wird bei \emph{Massendaten} praktisch schwierig; führt ggf. zu Bedarf an \emph{Checker-Tools} (die selbst qualifiziert werden müssten).
\end{itemize}

\paragraph{Kurzfazit.}
Aktuelle Luftfahrtstandards tragen die Grundidee von Assurance, stoßen bei ML aber an Grenzen: fehlende, anerkannte Verfahren für Generalisierung/Drift, Dissimilarität, ML-Coverage, probabilistische Leistungsnachweise, fortlaufende Updates und Tool-Einordnung.

\subsection{Limitierungen:}
Neben regulatorischen Lücken bestehen technologische Grenzen, die derzeit verhindern, System-Safety vollumfänglich zu garantieren. Die wichtigsten Punkte:

\paragraph{(1) Spezifikation.}
ML-Verhalten wird nur indirekt über Daten, Metriken und Lernverfahren beschrieben. ODD-Definitionen stoßen an Grenzen (Vollständigkeit, Terminologie, Annahmen, Edge-Cases); Off-Nominal-Bedingungen sind schwer abzudecken. Operative Daten können Spezifikationen nachschärfen, lösen aber selten Randfall-Probleme.

\paragraph{(2) Datenabhängigkeit \& Nachvollziehbarkeit.}
Leistung spiegelt Trainingsdaten (Bias, Lücken, Rauschen). Repräsentativität ist schwer sicherzustellen. Komplexe Pipelines (Preprocessing, Training, Konvertierung, Deployment) erschweren Traceability und Reproduzierbarkeit.

\paragraph{(3) Verifikation.}
Es fehlt ein allgemein akzeptiertes Maß für „ausreichendes Testen“. Formale Verfahren sind architektur-/aktivierungsabhängig und skaliert schwer. Struktur-Analogien (z.\,B. Neuron-Coverage) taugen nicht als SW-Coverage-Ersatz. XAI (lokal/global) hilft beim Verstehen/Debuggen, ersetzt aber keine belastbaren Test- und Robustheitsnachweise; Generalisierung und OoD bleiben offene Baustellen.

\paragraph{(4) Safety-getriebenes Lernen.}
Training optimiert meist Performance statt Sicherheit/Robustheit. Konfidenzwerte sind auf OoD-Daten unzuverlässig. Runtime-Assurance erfordert messbare Sicherheitskriterien und einfache, verifizierbare Fallback-Funktionen – bei komplexen Systemen schwer umzusetzen.

\paragraph{(5) Nichtdeterminismus.}
Variabilität entsteht durch Eingangsrauschen (Inference), parallelisierte Hardware (GPU/Multicore) und Trainingsstochasticität (Seeds, Shuffling). Erforderlich sind Monitoring/Logging, Sensitivitäts- und Robustheitsanalysen sowie Nachweise, dass Rest-Nichtdeterminismus vorhersehbar und akzeptabel begrenzt ist.

\paragraph{(6) Transfer Learning \& COTS.}
Transfer Learning reduziert Datenbedarf/Trainingszeit, birgt aber Risiken (negativer Transfer, Privatsphäre, Interpretierbarkeit). Fehlende Benchmarks behindern Vergleichbarkeit von Verifikationsmethoden. Für COTS-Modelle ist der bidirektionale Informationsfluss (Trainingsannahmen \(\leftrightarrow\) ODD) ungeklärt; unklar ist auch, wie nah der Prozess an klassischer Software-Integration liegt.

\paragraph{(7) Kontinuierliche Entwicklung (MLOps).}
ML benötigt laufendes Monitoring, Drift-Erkennung und kontrollierte Updates. Punktuelle Zertifizierungen passen schlecht zu häufigen Modelländerungen. Es fehlen klare Regeln, \emph{wann} retrainiert werden muss und \emph{welche} V\&V-Tiefe vor/nach Deployment erforderlich ist (Stichwort: kontinuierliche Re-Zertifizierung).

\paragraph{Kurzfazit.}
Ohne belastbare Verfahren für Spezifikation, Generalisierung/OoD, ML-Coverage, UQ, Runtime-Assurance, Traceability über die gesamte Pipeline und praktikable Re-Zertifizierungsprozesse bleibt „sicher genug“ für KI in der Luftfahrt eine offene Forschungs- und Standardisierungsaufgabe.







\clearpage
\section{Militärische Perspektive – Hauptunterschiede zur zivilen Luftfahrt}
In der Militärischen Luftfahrt findet bereits heute KI viele Anwendungen in den verschiedensten Bereichen. 
Im Hinblick auf so genannten \textit{``dual-use''} eröffnen sich viele Möglichkeiten bei denen sowohl die Zivile- 
als auch die Militärische- Branche profitieren können. Mit Fokus dieser Arbeit wird z.B. bereits im Eurofighter 
eine von Henshold entwickelte Objektklassifikations KI im Radar verwendet.
\textit{`` Standardization and certification on military AI applications ''} fasst die Hauptunterschiede zwischen Ziviler und Militärischer 
Luftfahrt wie folgt zusammen:\\
Militärische und zivile Verfahren unterscheiden sich grundlegend in Zielsetzung, Einsatzkontext, Regelwerk und Zertifizierung. 
Zivile Luftfahrt priorisiert Passagiersicherheit, Transparenz und internationale Harmonisierung; militärische Luftfahrt fokussiert 
Missionsfähigkeit, Resilienz und Geheimhaltung.

\paragraph{Vergleich in Kurzform.}
\begin{center}
\begin{tabular}{p{0.30\textwidth} p{0.30\textwidth} p{0.31\textwidth}}
\toprule
\textbf{Aspekt} & \textbf{Zivile Luftfahrt} & \textbf{Militärische Luftfahrt} \\
\midrule
Ziele \& Kontext & Sicherheit, Zuverlässigkeit, Interoperabilität im regulierten Linien-/Frachtbetrieb & Missionswirkung, Robustheit, Anpassungsfähigkeit unter widrigen Bedingungen \\
Regulatorik \& Standards & EASA/FAA/ICAO; DO-178C, DO-254, DO-160; KI-Guidance ergänzt bestehende Safety-Regeln & EDA/NATO/nationale Ministerien; EMARs, MIL-STD-882E; mögliche KI spezifische Ergänzungen \\
Standardisierung & Breite Stakeholder-Beteiligung, öffentlich, iterativ & Begrenzte Stakeholder (Behörden, Auftragnehmer), teils klassifiziert \\
Zertifizierung & Pflicht vor Einsatz; internationale Anerkennung; umfangreiche V\&V in kontrollierten Szenarien & EMAR-basiert und missionsbezogen; V\&V auch in Übungen/realitätsnahen Stresstests \\
Transparenz \& Ethik & Hohe Transparenz, Nachvollziehbarkeit, öffentliche Rechenschaft & Abwägung mit Geheimhaltung; zusätzliche ethische Fragen (z.\,B.\ autonome Wirkmittel) \\
Betrieb \& Resilienz & „Fail-safe“/„fail-operational“ nach Hazard; klarer Regelbetrieb & Stark auf „graceful degradation“, Wirksamkeit trotz Störung/Jamming/Täuschung \\
\bottomrule
\end{tabular}
\end{center}

\paragraph{Kernaussagen.}
\begin{itemize}
  \item \textbf{Ziele:} Zivil: Sicherheits- \& Compliance-getrieben = maximale Sicherheit, Zuverlässigkeit und Zertifizierbarkeit; Militär: Missions- \& Resilienz-getrieben = maximale Performance und operationelle Überlegenheit.
  \item \textbf{Regelwerke:} Zivil nutzt global harmonisierte, transparente Standards; Militär adaptiert und erweitert mit Fokus auf Einsatzanforderungen und Geheimschutz.
  \item \textbf{Nachweise:} Zivil zertifiziert vor Inbetriebnahme mit international anerkannten Nachweisen; Militär prüft zusätzlich unter extremen, unvorhersehbaren Bedingungen.
  \item \textbf{Ethik/Transparenz:} Zivil betont öffentliche Vertrauensbildung; Militär balanciert Transparenz mit operativer Sicherheit und gesonderten ethischen Fragestellungen.
\end{itemize}










\chapter{Automobilbranche: KI \& Sicherheit}

\section{Ziel \& Scope}
\textit{Zweck:} Rahmen für KI-Safety im Fahrzeug (Perzeption/Entscheidung, z.\,B.\ Radar-Objektklassifikation), ohne Security-Vertiefung. 
\textit{Hinweis:} Cybersecurity (ISO 21434, UNECE R155/R156) wird nur als Betriebsrandbedingung erwähnt und ist \emph{nicht} Teil dieser Thesis.
\textit{Gremien in einem Satz:} UNECE/WP.29 regelt Betrieb/Updates; ISO/SAE/IEEE standardisieren – Details hier nachrangig.

\subsection{Begriffe \& Standards – Schnellvergleich}
\begin{center}
\begin{tabular}{p{0.38\textwidth} p{0.25\textwidth} p{0.25\textwidth}}
\hline
\textbf{Aspekt} & \textbf{Luftfahrt} & \textbf{Automotive} \\
\hline
Safety-Level & DAL A–E & ASIL A–D \\
System-Safety & ARP4754A / ARP4761A & ISO 26262 (HARA, Safety Concept) / ISO 21448 (SOTIF) \\
Software-Assurance & DO-178C & ISO 26262 Teil SW \\
Hardware-Assurance & DO-254 & ISO 26262 Teil HW \\
KI-Guidance & EASA AI Level 1/2 (Guidance) & ISO/PAS 8800 \\
\hline
\end{tabular}
\end{center}

\section{Normen-Panorama}
\begin{itemize}
  \item \textbf{ISO 26262 (Functional Safety, ASIL A–D):} Lebenszyklus- und Nachweisrahmen für E/E-Systeme.
  \item \textbf{ISO 21448 (SOTIF):} Umgang mit Funktions-/Wahrnehmungs­insuffizienzen und Szenario-Lücken.
  \item \textbf{ISO 8800:} KI-spezifische Ergänzung zu 26262/SOTIF (Daten-/Modell-Governance, ODD/OOD, Kalibrierung, Betrieb/Updates).
  \item \textbf{ISO TR 5469:} Domänenübergreifender Leitfaden \emph{Funktionale Sicherheit \& KI} – Eigenschaften, Risiken, Methoden/Prozesse.
  \item \textbf{ISO 22989:} Domänenübergreifende Sammlung an KI Begriffen und Konzepten, die zur Vereinheitlichung dient.
\end{itemize}

\section{Deep Dive: ISO/PAS 8800}
\subsection{Positionierung \& Anwendungsbereich}
\textit{Kernidee:} Ergänzt 26262/SOTIF; kein „Assurance-Rabatt“ für KI; relevant bei ML in sicherheitskritischen Funktionen.

\subsection{Kernprinzipien (Stichworte)}
\begin{itemize}
  \item ODD/Use-Cases klar begrenzen; Unknown/Abstain + Fallback definieren.
  \item Daten-/Modell-Governance (Lineage, Versionierung, Label-Qualität, Reproduzierbarkeit).
  \item Unsicherheiten \& Kalibrierung; OOD-Erkennung als Pflichtbaustein.
  \item Szenario-basierte V\&V pro ODD-Segment (statt Mittelwert-KPIs).
  \item Betrieb/Monitoring \& kontrollierte Updates (inkl.\ Rollback).
\end{itemize}

\subsection{Artefakt-Checkliste für den Safety Case}
\textit{Entwürfe:} ODD-Spezifikation; Data/Model-Lineage; Kalibrierungs- \& OOD-Nachweise; Szenariomatrix \& Abdeckung; Runtime-Monitoring-Konzept; Change-/Rollback-Prozess.

\subsection{Beispiel: Radar-Objektklassifizierung}
\textit{Skizze:} Daten (RD/RDA, Mikro-Doppler, Tracks), Szenarien (Nacht/Regen/urban), Fehlklassifikationskosten, Abstain-Trigger, Drift-Signale, Integration in Brems-/Fusion-Kette, Fail-safe vs.\ Fail-operational.

\section{ISO/IEC TR 5469 – Brücke \& Ergänzung}
\textit{Rolle:} Referenzrahmen für KI+Funktionale Sicherheit; klärt Einsatzarten (KI \emph{in} Safety-Funktion, Safety \emph{für} KI, KI \emph{für} Entwicklung) und methodische Bausteine (z.\,B.\ Diversität/Unabhängigkeit von KI-Kanälen). 
\textit{Nutzen:} Ergänzt 8800/26262/SOTIF um Begriffe, Risikofaktoren und prozessuale Leitplanken – gut anschlussfähig an Luftfahrt-Argumentationsmuster.

\section{Crosswalk: Luftfahrt \(\leftrightarrow\) Automotive}
\subsection{Gemeinsamkeiten (kurz)}
\begin{itemize}
  \item Hazard-basierte Skalierung der Nachweisstrenge (DAL \(\leftrightarrow\) ASIL); kein „Assurance-Rabatt“ für KI.
  \item Safety Case/Assurance Case als roter Faden (Claims, Strategie, Evidenz, Annahmen).
  \item ODD/Operational Envelope als Einsatzgrenzen; Bedeutung von OOD/Unknown und Fallback/Degradation.
  \item Traceability, V\&V über Szenarien; Tool-Betrachtungen bei ML-Toolchains.
\end{itemize}

\subsection{Unterschiede (kurz)}
\begin{itemize}
  \item Regulatorik: Luftfahrt (Zulassung vor Betrieb, starke Behördenbindung) vs.\ Automotive (EU Typgenehmigung/UNECE, USA Self-Certification) – kürzere OTA-Zyklen im Feld.
  \item HMI \& Betrieb: Crew/Operator vs.\ Fahrer/Handover; höhere Variantenvielfalt, dichter Update-Betrieb im Automotive.
  \item Umgebungen: Luftfahrt eher kontrolliert/hochintegriert; Automotive breit gestreute ODDs (urban/land/\ldots).
\end{itemize}



\section{V\&V-Strategie (Automotive-konkret)}
\textit{Leitplanken:} Szenario-basiertes Testen je ODD-Segment; Kalibrierung \& OOD-Rate als Pflichtmetriken; KPI-Berichte pro Klasse/Szene/SNR/Range; Latenz- \& Degradationsregeln dokumentieren.

\section{Betrieb \& Lebenszyklus}
\textit{Kerngedanke:} Monitoring/Drift/Telemetrie festlegen; kontrollierte Updates mit Freigabe/Rollback; Pflege der Evidenz über Modellstände (Lineage). 
\textit{Hinweis:} Security/OTA-Regeln nur als Randbedingung, keine Vertiefung.

\section{Offene Punkte (normnah)}
\textit{To watch:} ML-Coverage/„genug getestet“, belastbare UQ-Nachweise, Transfer-Learning/COTS-Transparenz, Nichtdeterminismus, praktikable Prozesse für häufige Modell-Änderungen (kontinuierliche Re-Zertifizierung).
